% Part: incompleteness
% Chapter: incompleteness-provability
% Section: first-incompleteness-thm

\documentclass[../../../include/open-logic-section]{subfiles}

\begin{document}

% SEGMENT OLP-0315-S01

\olfileid{inc}{inp}{1in}

\olsection{முதல் முழுமையின்மைத் தேற்றம்}

முதல் முழுமையின்மைத் தேற்றத்துக்கான கோடலின் முதற் நிறுவலை இப்போது
விவரிக்கலாம். எண்கணித மொழியை நீட்டிக்கும் ஒரு மொழியில்
கணிக்கத்தக்கவாறு அடிகோளாக்கப்பட்ட எந்தக் கோட்பாடு~$\Th{T}$ ஐயும்
எடுத்துக்கொள்க; மேலும் $\Th{T}$, $\Th{Q}$ இன் அடிகோள்களை
உள்ளடக்கியதாக இருக்கட்டும். இதன் பொருள், குறிப்பாக, $\Th{T}$
கணிக்கத்தக்க சார்புகளையும் உறவுகளையும் பிரதிநிதித்துவப்படுத்துகிறது.

வாய்பாடுகளையும் நிறுவல்களையும் எண்களாக நியாயமான முறையில்
குறியீட்டாக்கினால், $\Prf[\Th{T}](x,y)$ என்ற உறவு கணிக்கத்தக்கது
என்று நாம் வாதிட்டுள்ளோம். இங்கு
$\Prf[\Th{T}](x,y)$ என்பது, $x$ ஆனது $\Th{T}$ இல்
$y$ என்ற கோடல் எண்ணுள்ள !!{formula} இன் !!a{derivation} இன்
கோடல் எண் என்பதையும் அப்போதும் அப்போது மட்டுமே குறிக்கிறது.
உண்மையில், கோடல் கருதிய குறிப்பிட்ட கோட்பாட்டுக்கு, தன்
கட்டுரையில் பட்டியலிட்ட 45 சார்புகளையும் உறவுகளையும் பயன்படுத்தி,
இந்த உறவு முதன்மை மீள்வரையறை என்பதை அவர் காட்ட முடிந்தது. அந்தப்
பட்டியலிலுள்ள 45 ஆம் உறவு, $x B y$, அவர் தேர்ந்தெடுத்த
$\Th{T}$ க்கான $\Prf[\Th{T}](x,y)$ என்பதே. கோடல் தன் கட்டுரையில்
“recursive” என்ற சொல்லைப் பயன்படுத்தும் இடங்களில், நாம் இன்று
“primitive recursive” என்ற சொற்றொடரைப் பயன்படுத்துவோம் என்பதை
நினைவில் கொள்க.

$\Prf[\Th{T}](x,y)$ கணிக்கத்தக்கது என்பதால், அது $\Th{T}$ இல்
பிரதிநிதித்துவப்படுத்தப்படுகிறது. அதைப் பிரதிநிதித்துவப்படுத்தும்
வாய்பாட்டைக் குறிக்க $\OPrf[\Th{T}](x,y)$ என்று எழுதுவோம். மேலும்
$\OProv[\Th{T}](y)$ என்பதை
$\lexists[x][\OPrf[\Th{T}](x,y)]$ என்ற வாய்பாடாக எடுத்துக்கொள்க.
இது கோடலின் பட்டியலிலுள்ள 46 ஆம் உறவான
$\fn{Bew}(y)$ ஐ விவரிக்கிறது. கோடல் குறிப்பிடுவது போல, இது
“recursive என்று கூற முடியாத” ஒரே உறவு. அவர் சொல்ல நினைத்தது
அநேகமாக இதுதான்: வரையறையிலிருந்து இது கணிக்கத்தக்கது என்பது
தெளிவாகத் தெரியவில்லை; பின்னர் ஏற்பட்ட வளர்ச்சிகள், உண்மையில்,
இது கணிக்கத்தக்கதல்ல என்பதை காட்டுகின்றன.

	$\Th{T}$ ஐக் கொண்ட !!{axiomatizable} கோட்பாடு ஒன்றை எடுத்துக்கொள்க;
	மேலும் அது~$\Th{Q}$ ஐக் கொண்டது.
அப்போது $\Prf[\Th{T}](x, y)$ தீர்மானிக்கத்தக்கது; எனவே அது
$\Th{Q}$ இல் !!a{formula}~$\OPrf[\Th{T}](x, y)$ மூலம்
பிரதிநிதித்துவப்படுத்தப்படுகிறது. மேலே விவரித்த வாய்பாட்டையே
$\OProv[\Th{T}](y)$ என்று எடுத்துக்கொள்க. நிலைப்புள்ளித்
துணைத்தேற்றத்தின் மூலம், $!G_\Th{T}$ என்ற வாய்பாடு ஒன்று உள்ளது;
அது $\Th{Q}$ (அதனால் $\Th{T}$ உம்) பின்வருவதை !!{derive}s
செய்யுமாறு அமைகிறது:
\begin{equation}
\ollabel{eqn:qpf}
!G_\Th{T} \liff \lnot \OProv[\Th{T}](\gn{!G_\Th{T}}).
\end{equation}
$!G_\Th{T}$ சாராம்சத்தில் “$!G_\Th{T}$ என்பது $\Th{T}$ இல்
!!{derivable} அல்ல” என்று கூறுகிறது என்பதை நினைவில் கொள்க.

\begin{lem}\ollabel{lem:cons-G-unprov}
$\Th{T}$ முரணற்ற, !!{axiomatizable} கோட்பாடாகவும்
$\Th{Q}$ ஐ நீட்டிப்பதாகவும் இருந்தால்,
$\Th{T} \Proves/ !G_\Th{T}$.
\end{lem}

\begin{proof}
$\Th{T}$, $!G_\Th{T}$ ஐ !!{derive}s செய்கிறது என்று கருதுக.
அப்போது !!a{derivation} ஒன்று \emph{இருக்கிறது}; எனவே ஏதோ ஒரு
எண்~$m$ க்கு $\Prf[\Th{T}](m,\Gn{!G_\Th{T}})$ என்ற உறவு
மெய்யாகும். ஆகவே $\Th{Q}$,
$\OPrf[\Th{T}](\num m, \gn{!G_\Th{T}})$ என்ற வாக்கியத்தை
!!{derive}s செய்யும். எனவே $\Th{Q}$,
$\lexists[x][\OPrf[\Th{T}](x,\gn{!G_\Th{T}})]$ ஐ
!!{derive}s செய்யும்; வரையறைப்படி இது
$\OProv[\Th{T}](\gn{!G_\Th{T}})$ ஆகும். \olref{eqn:qpf} இன்
மூலம் $\Th{Q}$, $\lnot !G_\Th{T}$ ஐ !!{derive}s செய்யும்;
$\Th{T}$ ஆனது $\Th{Q}$ ஐ நீட்டிப்பதால், $\Th{T}$ உம் அதை
செய்யும். ஆகவே $\Th{T}$, $!G_\Th{T}$ ஐ !!{derive}s
செய்தால், $\lnot !G_\Th{T}$ ஐயும் !!{derive}s செய்யும் என்று
காட்டியுள்ளோம்; எனவே அது முரணாகிவிடும்.
\end{proof}

\begin{defn}
\ollabel{thm:oconsis-q}
பின்வரும் நிபந்தனை பொருந்தினால் ஒரு கோட்பாடு~$\Th{T}$
\emph{$\omega$-முரணற்றது} எனப்படும்: எந்த
$\lexists[x][!A(x)]$ என்ற வாக்கியத்திற்கும், $\Th{T}$ ஆனது
$\lnot !A(\num 0)$, $\lnot !A(\num 1)$, $\lnot !A(\num 2)$,
\dots ஆகியவற்றை !!{derive}s செய்தால், $\Th{T}$ ஆனது
$\lexists[x][!A(x)]$ ஐ நிறுவாது.
\end{defn}

ஒவ்வொரு $\omega$-முரணற்ற கோட்பாடும் முரணற்றது என்பதை நினைவில்
கொள்க. இது எளிதாகப் பின்வருவதைக் கவனியுங்கள்: $\Th{T}$ முரணானது
என்றால், ஒவ்வொரு~$!A$ க்கும் $\Th{T} \Proves !A$. குறிப்பாக,
$\Th{T}$ முரணானால் ஒவ்வொரு~$n$ க்கும்
$\lnot !A(\num n)$ ஐயும், அதேபோல்
$\lexists[x][!A(x)]$ ஐயும் !!{derive}s செய்யும். ஆகவே $\Th{T}$
முரணானது எனில் அது $\omega$-முரணற்றது அல்ல. எதிர்மாற்றத்தின்
மூலம், $\Th{T}$ $\omega$-முரணற்றது என்றால் அது முரணற்றதாக
இருக்க வேண்டும்.

\begin{lem}\ollabel{lem:omega-cons-G-unref}
$\Th{T}$ ஒரு $\omega$-முரணற்ற, !!{axiomatizable} கோட்பாடாகவும்
$\Th{Q}$ ஐ நீட்டிப்பதாகவும் இருந்தால்,
$\Th{T} \Proves/ \lnot !G_\Th{T}$.
\end{lem}

\begin{proof}
$\Th{T}$, $\lnot !G_\Th{T}$ ஐ !!{derive}s செய்தால் அது
$\omega$-முரணற்றது அல்ல என்பதை காட்டுவோம். $\Th{T}$,
$\lnot !G_\Th{T}$ ஐ !!{derive}s செய்கிறது என்று கருதுக.
$\Th{T}$ முரணானது என்றால் அது $\omega$-முரணற்றது அல்ல; ஆகவே
முடிவு கிடைத்துவிட்டது. இல்லையெனில் $\Th{T}$ முரணற்றது. எனவே
\olref{lem:cons-G-unprov} இன் படி அது $!G_\Th{T}$ ஐ !!{derive}
செய்யாது. $\Th{T}$ இல் $!G_\Th{T}$ க்கான !!{derivation} எதுவும்
இல்லாததால், $\Th{Q}$ பின்வரும் வாக்கியங்களை !!{derive}s செய்யும்:
\[
\lnot \OPrf[\Th{T}](\num 0, \gn{!G_\Th{T}}), \lnot \OPrf[\Th{T}](\num 1,
\gn{!G_\Th{T}}), \lnot \OPrf[\Th{T}](\num 2, \gn{!G_\Th{T}}), \dots
\]
அதேபோல் $\Th{T}$ உம் அவற்றை !!{derive}s செய்யும். மறுபுறம்,
\olref{eqn:qpf} இன் படி $\lnot !G_\Th{T}$ என்பது
$\lexists[x][\OPrf[\Th{T}](x,\gn{!G_\Th{T}})]$ க்கு நிகரானது.
ஆகவே $\Th{T}$ $\omega$-முரணற்றது அல்ல.
\end{proof}

\begin{prob}
ஒவ்வொரு $\omega$-முரணற்ற கோட்பாடும் முரணற்றது. இதன் மறுதிசை
பொருந்தாது, அதாவது முரணற்றதாயினும் $\omega$-முரணற்றவை அல்லாத
கோட்பாடுகள் உள்ளன என்று காட்டுக. இதற்காக
$\Th{Q} \cup \{\lnot !G_\Th{Q}\}$ முரணற்றது, ஆனால்
$\omega$-முரணற்றது அல்ல என்பதை நிறுவுக.
\end{prob}

\begin{thm}
\ollabel{thm:first-incompleteness}
$\Th{T}$ ஐ, $\Th{Q}$ ஐ நீட்டிக்கும் எந்த $\omega$-முரணற்ற,
!!{axiomatizable} கோட்பாடான $\Th{T}$ உம் நிறைவானது அல்ல.
\end{thm}

\begin{proof}
$\Th{T}$ $\omega$-முரணற்றது என்றால் அது முரணற்றது; ஆகவே
$\Th{T} \Proves/ !G_\Th{T}$ என்பதை
\olref{lem:cons-G-unprov} காட்டுகிறது. மேலும்
\olref{lem:omega-cons-G-unref} இன் படி
$\Th{T} \Proves/ \lnot !G_\Th{T}$. இதன் பொருள் $\Th{T}$
நிறைவற்றது; ஏனெனில் அது $!G_\Th{T}$ ஐயும்
$\lnot !G_\Th{T}$ ஐயும் !!{derive}s செய்யவில்லை.
\end{proof}

\end{document}
